% Generated draft: requires mathematical and source review.
\documentclass[11pt]{article}
\usepackage[margin=25mm]{geometry}
\usepackage{fontspec}
\setmainfont{Latin Modern Roman}
\usepackage{amsmath,amssymb,mathtools,unicode-math}
\setmathfont{Latin Modern Math}
\usepackage{enumitem,needspace,xurl,hyperref}
\setlength{\parindent}{0pt}
\setlength{\parskip}{5pt}
\setlength{\emergencystretch}{3em}
\newcommand{\N}{\mathbb{N}}
\newcommand{\Z}{\mathbb{Z}}
\newcommand{\Q}{\mathbb{Q}}
\newcommand{\R}{\mathbb{R}}
\newcommand{\C}{\mathbb{C}}
\newcommand{\F}{\mathbb{F}}
\newcommand{\A}{\mathbb{A}}
\newcommand{\PP}{\mathbb{P}}
\newcommand{\E}{\mathbb{E}}
\newcommand{\eps}{\varepsilon}
\newcommand{\dd}{\,\mathrm{d}}
\begin{document}
\section*{UnsolvedMath ID: 701225}
\textbf{Source problem code:} catalogue entry
\subsection{Definitions and mathematical statement}
For a prime \(p\), a finite \(p\)-group is a finite group whose order is a power of \(p\). For a subgroup \(B\leq G\), write
\[
B^g=g^{-1}Bg,\qquad B^G=\{B^g:g\in G\},\qquad N_G(B)=\langle B^G\rangle.
\]
Here \(N_G(B)\) denotes the normal closure of \(B\) solely for this displayed definition; to avoid confusion with the usual normalizer notation, below we write the normal closure explicitly as \(\langle B^G\rangle\). Thus the required condition is
\[
B\trianglelefteq\langle B^G\rangle,
\]
meaning that \(x^{-1}Bx=B\) for every \(x\in\langle B^G\rangle\). An elementary abelian \(p\)-group is a group isomorphic to \((\Z/p\Z)^n\) for some integer \(n\geq0\).

The following three questions are attributed to G. Glauberman:
\begin{enumerate}
\item[(a)] For every integer \(n\geq0\) and every finite \(3\)-group \(G\), if \(G\) contains an abelian subgroup \(A\) of order \(3^n\), must there exist an abelian subgroup \(B\leq G\) such that
\[
|B|=3^n\quad\text{and}\quad B\trianglelefteq\langle B^G\rangle?
\]
\item[(b)] For every integer \(n\geq0\) and every finite \(2\)-group \(G\), if \(G\) contains an abelian subgroup \(A\) of order \(2^n\), must there exist an abelian subgroup \(B\leq G\) such that
\[
|B|=2^n\quad\text{and}\quad B\trianglelefteq\langle B^G\rangle?
\]
\item[(c)] For every integer \(n\geq0\) and every finite \(3\)-group \(G\), if \(G\) contains an elementary abelian subgroup \(A\) of order \(3^n\), must there exist an elementary abelian subgroup \(B\leq G\) such that
\[
|B|=3^n\quad\text{and}\quad B\trianglelefteq\langle B^G\rangle?
\]
\end{enumerate}
The supplied source reports that the corresponding assertions for primes \(p>3\) have been proved using extensions of Thompson's Replacement Theorem. It also reports that the dihedral group of order \(32\) gives a negative answer to the analogue of (c) with \(p=2\). These are supplied background claims, not independently verified conclusions of this draft. The local record labels the problem as open.
\subsection{Short English statement}
Given an abelian subgroup of a finite \(3\)-group or \(2\)-group, must there be an abelian subgroup of the same order that is normal in its own normal closure? For \(3\)-groups, ask also whether the replacement can be elementary abelian whenever the original subgroup is elementary abelian. The supplied record labels these questions as open.
\subsection{Sources}
\begin{itemize}
\item Supplied local UnsolvedMath registry record: permanent ID \(701225\), exact source problem code \texttt{catalogue entry}, title \emph{Kourovka Notebook Problem 16.33}. The record attributes the problem to G. Glauberman, proposed in 2006, and supplies the status \texttt{open}. Its \texttt{updated\_at} value is \texttt{2026-10-16T14:00:00Z}; this is registry metadata, not an independent literature-verification date.
\item \emph{The Kourovka Notebook (Unsolved Problems in Group Theory)}, Problem 16.33, 16th issue (2006); also identified in the supplied record with the 21st issue edition (2026), \texttt{arXiv:1401.0300v48}. The collection description names E. I. Khukhro and V. D. Mazurov as editors.
\item G. Glauberman, \emph{J. Algebra} \textbf{196} (1997), 301--338. Supplied as a reference for the background concerning greater primes and extensions of Thompson's Replacement Theorem.
\item J. Alperin and G. Glauberman, \emph{J. Algebra} \textbf{203} (1998), 533--566. Supplied as a reference for the same background.
\item G. Glauberman, \emph{J. Algebra} \textbf{272} (2004), 128--153. Supplied as a reference for the same background.
\end{itemize}
\end{document}
